\reading{CR-18}{Central Clearing}{DEFER}{1}

\begin{crkeybox}[Learning objectives for this reading]
\begin{enumerate}[label=\textbf{\alph*.},leftmargin=2em]
\item Define a central counterparty (CCP) and describe the mechanics of central clearing.
\item Explain the concept of novation under central clearing.
\item Define netting, multilateral offset, and compression and provide examples of each.
\item Describe the application and estimation of margin and default funds under central clearing.
\item Discuss the risks faced by a CCP and the ways it manages its exposures.
\item Provide examples of a loss waterfall.
\item Explain the different methods of absorbing losses and managing the default of one or more members of a CCP.
\item Compare bilateral and central clearing.
\item Compare initial margin and default fund requirements for clearing members in relation to loss coverage, cost of clearing, and moral hazard.
\item Describe the advantages and disadvantages of central clearing.
\end{enumerate}
\end{crkeybox}

\begin{crtrapbox}[Single-layer reading]
CR-18 exists only in the Crash layer.
\end{crtrapbox}

% =====================================================================
\lo{CR-18 a}{Define a central counterparty (CCP) and describe the mechanics of central clearing.}

CCPs act as intermediaries in OTC derivatives transactions, becoming the
\term{buyer to every seller and the seller to every buyer}.

\begin{itemize}
\item With a CCP there is increased \term{transparency}: CCPs have visibility into
all member positions, and counterparty risk is reduced.
\item Some parties might not access a CCP directly but through member firms,
adding complexity to margin transfers and default management.
\end{itemize}

\subsection*{Mechanics}

\begin{enumerate}[label=\alph*)]
\item CCPs earn revenue through fees on trades and income on held assets.
\item There are advantages and disadvantages to having few or many CCPs.
\textbf{Too many} CCPs drives competition up and credit quality down;
\textbf{too few} creates concentration risk.
\item \term{Criteria for clearing.} Standardization, complexity, liquidity,
wrong-way risk, and market volume are critical in determining whether a
derivative product can be cleared through a CCP.
\end{enumerate}

% =====================================================================
\lo{CR-18 b}{Explain the concept of novation under central clearing.}

\begin{crdefbox}[Novation]
When a party defaults, novation is the process by which the CCP replaces the
non-performing contract by stepping in as the new counterparty. The original
contractual obligations are \term{extinguished}, and the CCP assumes all
associated risks.
\end{crdefbox}

Once matched, transactions are guaranteed by the CCP, which strives to remain
\term{market neutral} by balancing buy and sell transactions.

% =====================================================================
\lo{CR-18 c}{Define netting, multilateral offset, and compression and provide examples of each.}

\begin{center}
\footnotesize
\begin{tabular}{L{3.0cm} L{11.5cm}}
\toprule
\thead{Term} & \thead{Definition} \\
\midrule
\term{Netting} & Reduces the number of transactions by combining or cancelling
out offsetting positions between \emph{two} parties, lowering counterparty risk \\
\term{Multilateral offset} & The same process applied across \emph{multiple}
parties, where the CCP centralizes and offsets the obligations, reducing overall
exposure among parties \\
\term{Compression} & Reducing the number of offsetting derivatives contracts
\emph{without impacting net exposure}. Cancels multiple contracts and replaces
them with fewer, lowering gross exposure \\
\bottomrule
\end{tabular}
\end{center}
\figcap{The worked compression example, with its net notional calculation, is at
\textbf{CR-16 e}, and the three-regime exposure comparison is at CR-16 c.}

% =====================================================================
\lo{CR-18 d}{Describe the application and estimation of margin and default funds under central clearing.}

CCPs use margins to reduce risk from member defaults.

\begin{enumerate}[label=\alph*)]
\item \term{Initial margin.} A required deposit to cover potential losses.
  \begin{itemize}
  \item The \emph{longer} it takes to unwind a defaulting member's positions, the
  \emph{higher} the initial margin requirement.
  \item The CCP aims to cover potential losses with near certainty, typically
  99\%, which again means higher initial margins.
  \end{itemize}
\item \term{Variation margin.} An additional deposit to reflect daily price
changes.
\end{enumerate}

% =====================================================================
\lo{CR-18 e}{Discuss the risks faced by a CCP and the ways it manages its exposures.}

\begin{enumerate}[label=\alph*)]
\item CCPs impose \term{membership requirements}, focused on creditworthiness,
liquidity, and likelihood of compliance with CCP rules.
\item \term{Margin} is the main form of risk mitigation in the event of member
default. Typically \emph{only cash} is accepted as variation margin, while cash
\emph{and highly liquid government securities} are accepted as initial margin.
\item \term{Member default} is the CCP's main risk. In a default the CCP reduces
its risk through \emph{macro-hedging} and \emph{auctions}.
\item The \term{margin period of risk (MPoR)} is the liquidation period. It can be
subdivided into three distinct periods: \textbf{pre-default}, \textbf{macro-hedging},
and \textbf{auctions}.
\end{enumerate}

% =====================================================================
\lo{CR-18 f}{Provide examples of a loss waterfall.}

\begin{center}
\begin{tikzpicture}[font=\sffamily\footnotesize,
  lay/.style={draw=FRMRule, line width=0.7pt, text width=6.4cm,
              align=center, minimum height=0.72cm}]
\node[lay, fill=PaleGreen, draw=FRMGreen] (l1) at (0,0)
  {\scriptsize Initial margin \emph{(defaulter)}};
\node[lay, fill=PaleGreen, draw=FRMGreen, below=1.5pt of l1] (l2)
  {\scriptsize Default fund \emph{(defaulter)}};
\node[lay, fill=PaleGold, draw=FRMGold, below=1.5pt of l2] (l3)
  {\scriptsize CCP skin-in-the-game};
\node[lay, fill=PaleRed, draw=FRMRed, below=1.5pt of l3] (l4)
  {\scriptsize Default fund \emph{(of non-defaulting members)}};
\node[lay, fill=PaleRed, draw=FRMRed, below=1.5pt of l4] (l5)
  {\scriptsize Rights of assessment and other loss allocation methods};
\node[lay, fill=PaleRed, draw=FRMRed, below=1.5pt of l5] (l6)
  {\scriptsize Remaining CCP capital};
\node[lay, fill=FRMSoft, draw=FRMGrey, below=1.5pt of l6] (l7)
  {\scriptsize Liquidity support, or the CCP fails};
% brackets
\draw[FRMGreen, line width=1pt] ($(l1.north east)+(0.25,0)$) --
  ($(l2.south east)+(0.25,0)$);
\node[right, font=\scriptsize\bfseries, text=FRMGreen, align=left]
  at ($(l1.east)!0.5!(l2.east)+(0.35,0)$) {DEFAULTER\\PAYS};
\draw[FRMRed, line width=1pt] ($(l4.north east)+(0.25,0)$) --
  ($(l7.south east)+(0.25,0)$);
\node[right, font=\scriptsize\bfseries, text=FRMRed, align=left]
  at ($(l4.east)!0.5!(l7.east)+(0.35,0)$) {SURVIVORS PAY\\\scriptsize\mdseries(moral hazard)};
\draw[FRMGold, dashed, line width=0.9pt] (-3.55,-1.98) -- (7.0,-1.98);
\draw[-{Stealth[length=6pt]}, FRMGrey, line width=1pt] (-4.1,0.35) -- (-4.1,-4.85)
  node[midway, left, font=\scriptsize\itshape, text=FRMGrey, rotate=90,
  anchor=south, yshift=4pt] {order of absorption};
\end{tikzpicture}
\end{center}
\figcap{A typical loss waterfall, defining the way in which the default of one or
more CCP members is absorbed. The boundary is the point worth memorising: the
\emph{CCP's own skin-in-the-game} sits between the defaulter's resources and the
surviving members' resources. Everything above it is paid by the party that
failed; everything below it is mutualised, which is where the moral hazard comes
from.}

% =====================================================================
\lo{CR-18 g}{Explain the different methods of absorbing losses and managing the default of one or more members of a CCP.}

\subsection*{Defaulter-pays resources}

\begin{itemize}
\item Initial margin and default fund contributions \emph{from the defaulter}.
\item Potential inclusion of additional variation margin held by the CCP.
\end{itemize}

\subsection*{Skin-in-the-game equity contribution}

The CCP's own equity contribution, to incentivize risk management and mitigate
losses beyond defaulter-pays resources.

\subsection*{Default fund}

A collective contribution from all clearing members. Loss allocation may be
\emph{pro rata}, or follow specific rules to encourage proactive participation in
auctions.

\subsection*{Additional absorption measures}

\begin{enumerate}
\item \term{Rights of assessment.} An \emph{unfunded} obligation to contribute
more if the default fund is significantly depleted, capped for risk control.
\item \term{Variation margin gains haircutting (VMGH).} Balancing losses by
reducing the \emph{gains} of other clearing members.
\item \term{Tear-up.} Terminating unmatched contracts to return to a matched book.
\item \term{Forced allocation.} Obliging members to accept certain portfolios at
CCP-determined prices.
\end{enumerate}

Beyond these lie the remaining CCP capital, and then liquidity support or failure
of the CCP.

% =====================================================================
\lo{CR-18 h}{Compare bilateral and central clearing.}

\begin{center}
\footnotesize
\begin{tabular}{L{2.6cm} L{5.7cm} L{6.2cm}}
\toprule
 & \thead{Bilateral} & \thead{Centrally cleared} \\
\midrule
Counterparty & Original & CCP \\
Products & All & Must be standard, vanilla, liquid \\
Participants & All & Clearing members are usually large dealers; other
  margin-posting entities can clear \emph{through} clearing members \\
Netting & Bilateral netting agreements and portfolio compression &
  Multilateral netting, including compression \\
Margining & Bilateral, bespoke arrangements dependent on credit quality;
  regulatory rules being introduced &
  Full collateralisation, including initial margin, enforced by the CCP \\
Close-out & Bilateral &
  Coordinated default management process (macro-hedges and auctions) \\
\term{Loss absorbency cost} & \textbf{Mainly capital} &
  \textbf{Mainly initial margin} (and default fund) \\
\bottomrule
\end{tabular}
\end{center}
\figcap{The last row is the one that reframes the whole comparison. Central
clearing does not remove the cost of bearing loss; it \emph{changes its form},
from capital held on a balance sheet to margin posted to a CCP.}

% =====================================================================
\lo{CR-18 i}{Compare initial margin and default fund requirements for clearing members in relation to loss coverage, cost of clearing, and moral hazard.}

\begin{center}
\footnotesize
\begin{tabular}{L{2.6cm} L{5.7cm} L{6.2cm}}
\toprule
 & \thead{Initial margin} & \thead{Default fund} \\
\midrule
Loss coverage & Provides high confidence coverage, e.g.\ 99\%, against member
  defaults &
  Being \emph{mutualized}, offers \textbf{higher} coverage than initial margin \\
Cost of clearing & Higher: each member funds its own &
  Makes clearing \textbf{cost-effective}, and therefore cheaper \\
Moral hazard & \textbf{Incentivizes better behaviour}: a member bears its own
  margin &
  \textbf{Increases moral hazard}: losses are shared \\
\bottomrule
\end{tabular}
\end{center}
\figcap{A clean three-way trade-off, and every row runs the same direction. The
default fund buys coverage and cheapness by mutualising, and mutualising is
exactly what weakens the incentive to behave.}

% =====================================================================
\lo{CR-18 j}{Describe the advantages and disadvantages of central clearing.}

\begin{center}
\footnotesize
\begin{tabular}{L{7.0cm} L{7.5cm}}
\toprule
\thead{Advantages} & \thead{Disadvantages} \\
\midrule
Transparency & Moral hazard \\
Multilateral offset & Adverse selection \\
Loss mutualization & Bifurcations \\
Legal and operational efficiency & Procyclicality \\
Liquidity enhancement & \\
Default management & \\
\bottomrule
\end{tabular}
\end{center}

\begin{crtrapbox}[Loss mutualization is on both lists]
Loss mutualization appears as an \emph{advantage} here and moral hazard as a
\emph{disadvantage}, and objective i makes the connection explicit: the
mutualisation that delivers coverage is the same mechanism that creates the moral
hazard. The fuller advantage and disadvantage lists in Gregory's chapter 2
framing are at \textbf{CR-14 g}.
\end{crtrapbox}
